\subsection{Short even multipairs with a nonunit ordinary anchor}\label{proof:R2-C038}
\paragraph{R38.1 Actual objects and directions}

Let the original receiving arm have length \(m\ge3\), and let the terminal's strict residual list be
\fitdisplay{
a_1<b_1<\cdots<a_j<b_j,\qquad j\ge2,\qquad a_1\in\{2,3\},\qquad
F=\prod_{i=1}^j h_{a_i}f_{b_i}.
}{R38.L74}
An optional negative root-post block is
\fitdisplay{
P=\prod_{i=1}^k W_{u_i}\bar f_{v_i},\qquad
4\le u_1<v_1<\cdots<u_k<v_k,\qquad k\ge0.
}{R38.L76}
All counts and lengths are finite. The two masks may overlap each other and the anchor length. Genuine equal-pair products \(E_t,\allowbreak E_r\in[0,\allowbreak 1]\) remain at their original vertices. Define the decorated message by inserting \(E_tF\) in the following connector and multiplying by \(E_rP\) at the receiving root:
\fitdisplay{
H_D=E_rP\begin{cases}
\bar f_mA(E_tF),&a_1=2,\ L=1,\\
W_mT^L(E_tF),&a_1=2,\ L\ge2,\\
W_mT^L(E_tF),&a_1=3,\ m\le L,\\
\bar f_mAT^{L-1}(E_tF),&a_1=3,\ m>L.
\end{cases}
}{R38.L78}
Every \(f\) arm points outward, every \(\bar f\) arm inward. An \(h\) arm has its first edge toward the root and all later edges toward the leaf; \(W\) reverses that entire arm. A \(T\) connector points toward the receiving root. In an \(AT^{L-1}\) row, only its root-adjacent edge is \(A\). These choices depend only on the original integers, and satisfy
\fitdisplay{
\boxed{EH_D^8\le e^{-2d},\qquad d=\|K\one\|_2^2.}
}{R38.L87}

\paragraph{R38.2 Finite inputs and one actual terminal mask}

Use the finite packet of R2-C025 and the unequal-pair envelopes of Q5-C090. The first pair has loss \(\kappa d\) at minimum two and \((2-\theta)d=51d/32\) at minimum three; every later shorter arm has length at least four and pays \(16d/5\). Thus
\fitdisplay{
F\le e^{R-\lambda d},\qquad
\lambda_2=607/160,\quad\lambda_3=767/160,\qquad
R=\sum_i\sum_{r=a_i+1}^{b_i}K^r\one.
}{R38.L104}
Extra pairs only enlarge the available loss. The intervals within this single finite mask are disjoint. Work throughout on the original measure, with
\fitdisplay{
\begin{gathered}
B=-K,\quad u=B\one,\quad T=J+B,\quad A=J-B,\quad p=1-u,\quad q=1+u,\\
\delta_i=\|B^i\one\|^2,\quad d=\delta_1,\quad
\theta=13/32,\quad\kappa=19/32,\quad\tau=d-\delta_2,\\
f_n=A^n\one,\quad\bar f_n=T^n\one,\quad h_n=Tf_{n-1},\quad W_n=A\bar f_{n-1}.
\end{gathered}
}{R38.L106}
The finite calculus gives \(0\le d\le1/3\), \(d^2\le\delta_2\le\theta d\), \(\delta_{i+1}\le\theta\delta_i\), \(\|B\|^2\le\theta\), and \(|BV(x)|\le\|V\|\). Only \(f_n,\allowbreak \bar f_n\) have means at most one; \(Eh_2=EW_2=1+d\). All four path types lie in \([0,\allowbreak 2]\). The common zero-boundary argument handles \(d=0\); below take \(d>0\).

At \(z=i\sqrt x\), let \(\nu\) be the original spectral measure of \(-B^2\) at \(u\), with mass \(d\) and \(\int x^i\,\allowbreak d\nu=\delta_{i+1}\). If \(R=U(B)u\), \(U=\alpha+z\beta\), skew adjunction gives
\fitdisplay{
ER=-\int\beta\,d\nu,\qquad
\|R\|^2=\int(\alpha^2+x\beta^2)\,d\nu.
}{R38.L115}
All masks are finite; positive geometric tails only bound them. Set
\fitdisplay{
G_0=0,\quad G_n(z)=\sum_{r=0}^{n-1}z^r,\quad\beta_1=\beta_2=0,\quad
\beta_n=\sum_{i=1}^{\lfloor(n-1)/2\rfloor}(-1)^{i-1}\delta_i.
}{R38.L124}
The exact finite positive convolution, equivalently \(S_n^-=\sum_{r=0}^nB^r\one=1+G_n(B)u\), is
\fitdisplay{
1+G_n(B)u=\bar f_n+\sum_{r=3}^n\beta_r\bar f_{n-r},\qquad
\tau\le\beta_r\le d\quad(r\ge3).
}{R38.L126}
Four alternating residue classes give \(|Q(i\sqrt x)|^2\le(1+x)/(1-x^2)^2\le32768/16245\) for every finite zero-one mask \(Q\). This supplies the norms below. At minimum two the mandatory term is \(K^3\one\). The short pair retains the actual \(\beta_bf_0\) term; \(\beta_2=0\) is never charged as a positive loss.

For minimum two, the full actual packet is
\fitdisplay{
\begin{gathered}
\sigma=ER^2<d/3,\quad |R|<3/5,\quad e^R\le1+R+5R^2/8,\\
R=U(B)u,\quad U=\alpha+z\beta,\quad
x/2\le\alpha\le6x/5,\quad-6x/5\le\beta\le x/2,\\
D_i=EB^iR,\quad a_0=-D_1\in[0,d/2],\quad
\rho=-D_2\in[-d/12,d/5],\\
D_3\in[0,d/5],\quad |D_i|\le d/12\ (i\ge4),\quad D_0\in[-d/5,d/2],\\
D=D_0+5\sigma/4\in[-d/5,11d/12],\qquad
\epsilon=D+D_1\in[-7d/10,7d/10].
\end{gathered}
}{R38.L134}
These are correlated quantities from the same \(R\). Positive propagation gives
\fitdisplay{
\begin{gathered}
e^{\lambda d}T^LF\le Y_L,\quad Y_1=q+D+BR,\\
Y_L=\bar f_L+B^LR+D\bar f_{L-1}
 +\sum_{i=1}^{L-1}D_i\bar f_{L-1-i},\\
Q_n=G_n+z^nU+\epsilon G_{n-2}+Dz^{n-2}\quad(n\ge2),\\
Q_2=G_2+D+z^2U,\qquad Q_{n+1}=1+\epsilon+zQ_n,\\
Y_L=1+\epsilon+Q_L(B)u-\sum_{n=3}^L\gamma_n\bar f_{L-n}.
\end{gathered}
}{R38.L142}
The exact signed coefficients satisfy
\fitdisplay{
\begin{gathered}
\gamma_3=d-D_2\ge11d/12,\qquad
\gamma_4=d+Dd-D_3\ge11d/15,\\
\gamma_n=\beta_n+D\beta_{n-1}-a_0\beta_{n-2}-D_{n-1}\ge d/4
\quad(n\ge5).
\end{gathered}
}{R38.L148}
The long bound uses \(\gamma_n/d\ge\kappa-1/12-(7/10)d\ge133/480>1/4\). The short \(\gamma_4\) bound retains \(D\ge-d/5\), \(D_3\le d/5\), \(d\le1/3\). Moreover \(EQ_n(B)u=-\gamma_{n+1}\) for all \(n\ge2\). At every \(L\ge4\), retaining the last two actual losses with \(h=1/4\) yields
\fitdisplay{
Y_L\le(1+\epsilon-hd)q-hd+BQ_{L-1}(B)u.
}{R38.L154}
Every discarded loss has its proved nonnegative sign at every finite index.

Positive convolution followed by \(A\) gives, for every \(m\ge3\),
\fitdisplay{
W_m\le p-\beta_m-BG_{m-1}(B)u,\qquad\beta_m\in[\kappa d,d].
}{R38.L158}
The actual increment \(\Delta_m=\beta_{m+1}-\beta_m\le\delta_3<d/6\): odd \(m\) gives zero, \(m=4\) gives \(-\delta_2\), and positive even increments first occur at \(m=6\), with \(\delta_3\) or later \(\delta_5,\allowbreak \ldots\). Every \(G_n\) contracts except \(G_2\), whose squared bound is \(1+\theta\), with \(\sqrt{1+\theta}<6/5\). The finite difference identity \(G_n-G_r=\pm z^{\min(n,\allowbreak r)}G_{|n-r|}\), or zero, gives
\fitdisplay{
\|(G_n-G_r)u\|<\sqrt d/2\quad(n,r\ge2),
\qquad\theta\sqrt{1+\theta}<1/2.
}{R38.L163}

\paragraph{R38.3 Full-mean eighth-exponent payment}

If the actual anchor plus its nonnegative propagated upper is at most \(2+c+BV\), then AM--GM and the all-real signed exponential lemma give
\fitdisplay{
\log EH^8\le-8\lambda d+8(c+EBV)+17\|V\|^2,
\qquad Ee^{8BV}\le e^{8EBV+17\|V\|^2}.
}{R38.L169}
The second inequality is the full fourth-exponent lemma applied to \(2V\). The actual nonnegative average also satisfies \(y^{16}\le e^{16(y-1)}\). The precursor and its signed mean are retained. A root-post block adds \(-8\Gamma_r+8ES\) and the single vector \(V_p\), as paid in \mbox{\hyperref[part:R38.8]{part R38.8}}.

\paragraph{R38.4 Minimum two, unit connector}

Use \(\bar f_mAF\). The positive Taylor error and actual mean are
\fitdisplay{
S_1R=AR,\quad\omega_1=5\sigma/4<5d/12,\quad
EAR=D_0-D_1\le d,
}{R38.L180}
using \(D_0,\allowbreak a_0\le d/2\). The bare sum has \(\bar f_m+p\le2-\beta_m+BV_b\), \(V_b=G_{m-1}u\), \(EBV_b=-\beta_{m+1}\). At \(m=3\), \(\beta_3+\beta_4=2d\); at \(m\ge4\) that sum is at least \(2\tau\ge19d/16\).

For \(m=3\),
\fitdisplay{
\langle G_2u,R\rangle=\int(\alpha+x\beta)\,d\nu\ge0,
\qquad\alpha+x\beta\ge x-\frac{2x^2}{1-x^2}\ge0.
}{R38.L187}
Here \(2\theta<1-\theta^2\) is the rational comparison \(832<855\) in units of \(1024\). Hence \(\|G_2u-R\|^2\le(45/32+1/3)d=167d/96\). For \(m\ge4\), \(\operatorname{Re}G\ge\kappa\), the normalized imaginary coefficient has modulus at most one, and the finite coordinate bounds imply
\fitdisplay{
\langle G_{m-1}u,R\rangle\ge-d/12.
}{R38.L191}
Indeed the concave function \(\kappa x/2-(6/5)x^2\) has its minimum at zero or \(\theta\), and its latter value is \(-793/10240>-1/12\); \(\beta\ge0\) only improves the dot product. Thus \(\|V_b-R\|^2\le d+d/3+d/6=3d/2\). The two complete loss payments are
\fitdisplay{
\begin{gathered}
607/20+16-17(167/96)-10/3-8>14/3>2\quad(m=3),\\
607/20+19/2-51/2-10/3-8=181/60>2\quad(m\ge4).
\end{gathered}
}{R38.L197}
For the first, \(607/20>30\) and \(17(167/96)<30\). All \(j\ge2\) are already included in the one finite mask.

\paragraph{R38.5 Minimum two, length-two connector}

Keep the actual \(W_mT^2F\) direction. From the genuine \(Y_2\),
\fitdisplay{
\begin{gathered}
W_m+Y_2\le2+c+BV,\quad c=\epsilon-\beta_m,\\
V=D\one-B(G_{m-2}u-R),\quad V'=B(G_{m-2}u-R),\\
EBV'=d-\beta_{m+1}-D_2,\\
c+EBV=\epsilon-d+\Delta_m+D_2\le\epsilon-3d/4\le-d/20.
\end{gathered}
}{R38.L204}
For \(n=m-2\ge1\), including \(n=1,\allowbreak 2\), each \(G_n\) has real coordinate in \([\kappa,\allowbreak 1]\) and normalized imaginary coordinate in \([0,\allowbreak 1]\), whereas \(\alpha\le6\theta/5=39/80<\kappa\). Thus
\fitdisplay{
|G_n-U|^2\le(1-x/2)^2+x(1+6x/5)^2
=1+(53/20)x^2+(36/25)x^3<8/5.
}{R38.L213}
This uses \(\theta^2<1/6\), \(\theta^3<7/100\), and \(53/120+63/625<3/5\). Hence \(\|V'\|^2\le13d/20\), while \(EV'=-\beta_m+a_0\) and \(\beta_m-a_0\ge3d/32>0\). Retain the full constant cross:
\fitdisplay{
\|V\|^2=\|V'\|^2+D^2+2D(\beta_m-a_0).
}{R38.L220}
If \(D<0\), the cross is nonpositive and \(D^2\le d/75\). If \(D\ge0\), use \(D\le11d/12\), \(\beta_m-a_0\le d\), \(d\le1/3\), so the extra cost is at most \(385d/432\). Consequently
\fitdisplay{
\begin{gathered}
\|V\|^2\le(13/20+385/432)d=3329d/2160<25d/16,\\
607/20+2/5-17(3329/2160)>17/4>2,
\end{gathered}
}{R38.L225}
where \(17(3329/2160)<53/2\).

\paragraph{R38.6 Minimum two, every longer connector}

At \(L=3\),
\fitdisplay{
\begin{gathered}
Y_3=(1+\epsilon)q+BQ_2u-\gamma_3,\quad c=\epsilon-\beta_m-\gamma_3,\\
V=V_0+\epsilon\one,\qquad
V_0=(G_2-G_{m-1})u+Du+B^2R,\\
\|V_0\|<(1/2+11/36+6/25)\sqrt d<21\sqrt d/20.
\end{gathered}
}{R38.L233}
The \(B^2R\) term uses \(\theta^2/3<36/625\). The mean \(EV_0=\beta_m-\gamma_3\in[-2d/3,\allowbreak d/12]\), with \(\gamma_3\in[11d/12,\allowbreak 6d/5]\). For \(\epsilon\ge0\), its constant correction is at most \(((7/10)^2+(7/10)/6)d^2\le91d/450\); for \(\epsilon\le0\), it is at most \(((7/10)^2+(4/3)(7/10))d^2\le427d/900\). Therefore
\fitdisplay{
\begin{gathered}
\|V\|^2<[(21/20)^2+427/900]d<169d/100,\\
c+EBV=\epsilon-\gamma_3-\gamma_4+\Delta_m
\le\epsilon-89d/60\le-47d/60,\\
607/20+94/15-17(169/100)>457/60>2.
\end{gathered}
}{R38.L243}
The last comparison uses \(17(169/100)<29\).

For every \(L\ge4\), put \(n=L-1\ge3\), \(h=1/4\), and retain
\fitdisplay{
c=\epsilon-2hd-\beta_m,\quad
V=V_0+(\epsilon-hd)\one,\quad V_0=(Q_n-G_{m-1})u.
}{R38.L251}
The four terms in \(V_0\) are the geometric difference, \(B^nR\), \(\epsilon G_{n-2}u\), and \(DB^{n-2}u\). Their respective bounds are
\fitdisplay{
\begin{gathered}
\theta\sqrt{1+\theta}<1/2,\quad \theta^3/3<9/400,\\
(7/10)(1/3)(6/5)=7/25,\quad
(11/12)(1/3)(13/20)=143/720<1/5.
\end{gathered}
}{R38.L256}
The actual \(G_2\) exception is paid by \(6/5\). Hence \(\|V_0\|<113\sqrt d/100\). Also \(EV_0=\beta_m-\gamma_L\in[-d,\allowbreak 3d/4]\): at \(L=4\), \(\gamma_4\le d+11d^2/12\) because \(D_3\ge0\); at \(L\ge5\), \(\gamma_L\le13d/12+11d^2/12\le25d/18\). Combining with \(\beta_m\ge\kappa d\) gives the lower mean bound, and \(\gamma_L\ge hd\) gives the upper. The signed exponent mean is
\fitdisplay{
c+EBV=\epsilon-2hd-\gamma_{L+1}+\Delta_m\le\epsilon-7d/12.
}{R38.L263}
Put \(t=\epsilon/d-1/4\in[-19/20,\allowbreak 9/20]\) and \(C=113/100\). The full squared norm obeys
\fitdisplay{
\frac{\|V\|^2}{d}\le N(t)=\begin{cases}
C^2+(t^2-2t)/3,&t\le0,\\
C^2+(t^2+3t/2)/3,&t\ge0.
\end{cases}
}{R38.L270}
The fee \(8(t-1/3)+17N(t)\) is convex on each closed half. At the negative, zero and positive endpoints, the pairs \((N,\allowbreak \text{mean fee})\) are
\fitdisplay{
(C^2+1121/1200,-154/15),\quad(C^2,-8/3),\quad(C^2+117/400,14/15).
}{R38.L276}
Every fee is less than \(607/20-2\). Indeed \(17C^2=217073/10000<87/4\); the positive endpoint uses \(17(117/400)=1989/400<5\), \(14/15<1\), giving fee \(<111/4\) and loss \(>13/5\). The negative endpoint uses \(17(1121/1200)<16\), \(154/15>41/4\), giving fee \(<55/2\) and loss \(>57/20\). At zero the fee is \(<87/4-8/3<20\), giving loss \(>207/20\). Convexity covers every allowable correlated \((\epsilon,\allowbreak D)\). This proves all finite \(L\ge4\) for minimum two in the same \(W_mT^L\) direction.

\paragraph{R38.7 Minimum three, every connector length}\label{part:R38.7}

The actual R2-C025 packet gives
\fitdisplay{
\begin{gathered}
\sigma<9d/64,\quad |R|<1/3,\quad e^R\le1+R+3R^2/5,\\
D_0\in[-d/5,d/2],\quad D_1\in[-d/12,d/5],\quad-D_2\in[-d/12,d/5].
\end{gathered}
}{R38.L293}
with \(\lambda=767/160\). Keep the structural comparison \(m\le L\) versus \(m>L\). The ordinary power bounds are \(\|T^n\|\le1\) for \(n\ge3\) and
\fitdisplay{
\|T^2\|\le(1+1024\delta_2/855)^2
\le(2981/2565)^2<(7/6)^2<7/5.
}{R38.L297}
The genuine positive Taylor errors have the exact scopes
\fitdisplay{
\omega_1=(6/5)\sigma<27d/160,\qquad
\omega_L<27d/80\quad(L\ge2).
}{R38.L301}
Moreover \(EAR\le7d/12\), \(EATR=D_0+D_1-D_2\le9d/10\), and
\fitdisplay{
\|TR\|^2=D_0^2+2D_0D_1+\|BR\|^2
\le(3/20+117/2048)d<(11/24)^2d.
}{R38.L302}
For \(L\ge3\), the power bounds give \(\|T^{L-1}R\|<21\sqrt d/40\). The finite Duhamel expansion
\fitdisplay{
f_k=\mu_{k-1}\one+\sum_{r=1}^{k-1}\mu_{k-1-r}(-B)^r\one+(-B)^k\one
}{R38.L310}
has coefficients in \([0,\allowbreak 1]\), whence \(\|f_k-\mu_{k-1}\one\|\le\sqrt d\sum_{r=0}^{k-1}\theta^{r/2}<3\sqrt d\). It gives \(a_R=Ef_{L-1}R\le d/2+3\sqrt{d\sigma}<13d/8\), and thus, for either actual last-edge sign,
\fitdisplay{
ES_LR\le13d/8+21d/40=43d/20\quad(L\ge3).
}{R38.L305}
Positivity and the last row bounded by two give \(S_LF\le e^{-\lambda d}(S_L\one+S_LR+\omega_L)\), nonnegative as a whole. At \(L\ge2\), the error uses \(E[f_{L-1}R^2]\le2\sigma\); at \(L=1\), it uses \(\sigma\). Retain \(S_LR=a_R\pm BT^{L-1}R\), including its true constant and the displayed constant cross in \(TR\).

The exact bare sums are:
\fitdisplay{
\begin{gathered}
\bar f_m+p\le2-\beta_m+BG_{m-1}u\quad(L=1),\\
8(\beta_m+\beta_{m+1})=16d\ (m=3),\qquad
8(\beta_m+\beta_{m+1})\ge19d/2\ (m\ge4),\\
\bar f_m+W_2\le2-\beta_m+BV_b,\quad
V_b=BG_{m-2}u,\quad\|V_b\|\le4\sqrt d/5,\\
-8\beta_m+8EBV_b\le-3d/2\quad(L=2),\\
W_h+\bar f_H\le2-(51/32)d+BV_b,\quad
EBV_b\le\theta d,\quad\|V_b\|\le\sqrt d/2\quad(L\ge3),\\
h=\min(m,L)\ge3,\qquad H=\max(m,L).
\end{gathered}
}{R38.L321}
For \(h=H\), the exact sum is \(2\mu_{h-1}\le2-2d\); otherwise the finite vector is \(B^{h-1}G_{H-h}u\). The \(G_2\) exception uses \(\theta^2(1+\theta)<1/4\), and the signed mean uses \(\delta_2\le\theta d\).

The four complete payments, in order \((L=1,\allowbreak m=3)\), \((L=1,\allowbreak m\ge4)\), \(L=2\), \(L\ge3\), are
\fitdisplay{
\begin{gathered}
\|G_2u-R\|<(6/5+3/8)\sqrt d=63\sqrt d/40,\\
767/20+16-17(63/40)^2-27/20-14/3>2,\\
\|G_{m-1}u-R\|<11\sqrt d/8,\quad
767/20+19/2-17(11/8)^2-27/20-14/3>2,\\
\|V\|\le(4/5+11/24)\sqrt d=(151/120)\sqrt d,\\
767/20+3/2-17(151/120)^2-27/10-36/5>2,\\
\|V\|\le(1/2+21/40)\sqrt d=(41/40)\sqrt d,\\
767/20+19/2-17(41/40)^2-27/10-86/5>2.
\end{gathered}
}{R38.L333}
Their norm costs are less than \(43,\allowbreak 33,\allowbreak 27,\allowbreak 18\), and their other fees are less than \(7,\allowbreak 7,\allowbreak 10,\allowbreak 20\). The resulting weaker losses are \(87/20,\allowbreak 157/20,\allowbreak 57/20,\allowbreak 197/20\), each greater than two. Equality \(m=L\) and every finite connector are included.

\paragraph{R38.8 Every finite negative root-post block}\label{part:R38.8}

For \(k\ge1\), the strong actual pair envelope and its one finite reversed mask give
\fitdisplay{
\begin{gathered}
P\le e^{S-\Gamma_r},\quad\Gamma_r\ge16kd/5,\quad
S=BV_p=B^4Q_p(B)u,\quad V_p=B^3Q_p(B)u,\\
C_\star=32768/16245,\quad C_\beta=1024/855,\\
\|V_p\|^2\le C_\star\theta^3d=2197d/16245<9d/64,\qquad
ES\le C_\beta\delta_4<d/12.
\end{gathered}
}{R38.L347}
Add \(V_p\) to the same actual \(V\) and pay every cross by the triangle inequality. Its intervals may overlap the terminal intervals and \(m\).

For minimum two, the short rows use base norm bounds \(4/3,\allowbreak 5/4,\allowbreak 13/10\), in units of \(\sqrt d\), and combined bounds \(41/24,\allowbreak 13/8,\allowbreak 67/40\). The full losses are
\fitdisplay{
\begin{gathered}
607/20+128/5+19/2-17(41/24)^2-(10/3+8+2/3)>69/20>2,\\
607/20+128/5+2/5-17(13/8)^2-2/3>207/20>2,\\
607/20+128/5+94/15-17(67/40)^2-2/3>793/60>2.
\end{gathered}
}{R38.L357}
The respective norm costs are below \(50,\allowbreak 45,\allowbreak 48\); the first other fees total twelve, and the last two use \(2/3<1\).

At minimum two and \(L\ge4\), complete the squares:
\fitdisplay{
N=C^2-1/3+(t-1)^2/3\quad(t\le0),\qquad
N=C^2-3/16+(t+3/4)^2/3\quad(t\ge0).
}{R38.L367}
Both residual constants are positive, so \(\sqrt N\) and \((\sqrt N+3/8)^2\) are convex on each half. Hence so is the full fee \(8(t-1/3)+17(\sqrt N+3/8)^2+2/3\). At the negative, positive and zero endpoints, respectively, the combined norms are less than \(15/8,\allowbreak 327/200,\allowbreak 301/200\), using \(\sqrt N<3/2\) and \(\sqrt N<63/50\) at the first two. Their fees are bounded by
\fitdisplay{
\begin{gathered}
17(15/8)^2-154/15+2/3,\\
17(327/200)^2+14/15+2/3,\\
17(301/200)^2-8/3+2/3.
\end{gathered}
}{R38.L375}
The norm costs are below \(60,\allowbreak 46,\allowbreak 39\), so the complete fees are below \(252/5,\allowbreak 238/5,\allowbreak 37\). Subtracting from \(607/20+128/5=1119/20\) leaves losses greater than \(111/20,\allowbreak 167/20,\allowbreak 379/20\), all greater than two. Composition with the increasing convex square on nonnegative arguments justifies the convexity and covers every actual \(\epsilon\).

For minimum three, the four combined norms are \(39/20,\allowbreak 7/4,\allowbreak 49/30,\allowbreak 7/5\), in the same order as \mbox{\hyperref[part:R38.7]{part R38.7}}. The complete payments are
\fitdisplay{
\begin{gathered}
767/20+128/5+16-17(39/20)^2-27/20-14/3-2/3>2,\\
767/20+128/5+19/2-17(7/4)^2-27/20-14/3-2/3>2,\\
767/20+128/5+3/2-17(49/30)^2-27/10-36/5-2/3>2,\\
767/20+128/5+19/2-17(7/5)^2-27/10-86/5-2/3>2.
\end{gathered}
}{R38.L385}
Their norm costs are below \(65,\allowbreak 53,\allowbreak 46,\allowbreak 34\), and the remaining fees below \(7,\allowbreak 7,\allowbreak 11,\allowbreak 21\). The starting losses and resulting net losses satisfy
\fitdisplay{
\begin{gathered}
\text{starting losses: }1599/20,\ 1469/20,\ 1309/20,\ 1469/20,\\
\text{net losses: }>159/20,\ >269/20,\ >169/20,\ >369/20.
\end{gathered}
}{R38.L393}
Extra terminal or root pairs increase the available loss. Restore the actual \(E_r,\allowbreak E_t\le1\) by positive operators to finish the theorem.
